\documentclass[12pt]{article}
\usepackage[utf8]{inputenc}
\usepackage[T1]{fontenc}
\usepackage{hyperref}
\title{Fagauvea Language Resources for K-12 Curriculum Development}
\author{Compiled by Claude Sonnet 5 (claude-sonnet-5) for Andr\'as Kornai}
\date{}
\begin{document}
\maketitle

\section{Introduction}
Fagauvea (West Uvean, ISO 639-3 uve) is a Polynesian outlier language spoken on Ouv\'ea in the Loyalty Islands of New Caledonia, and in Noum\'ea. The 2009 census gave about 2,200 speakers. It is closest to Wallisian, but about three centuries of contact with Iaai have changed its phonology and lexicon. It can be written in the Latin alphabet, with long vowels doubled or marked with a macron.

\section{Primary Educational Resources}
No primers or school textbooks for Fagauvea turned up in the searches.

\section{Linguistic and Reference Materials}
The search results list the following works only as bibliographic entries. Their availability online was not confirmed.
\begin{itemize}
\item K. J. Hollyman, \emph{De muna fagauvea I: dictionnaire fagauvea-fran\c{c}ais} (1987), and an earlier Emprin dictionary, both cited on the CRDO lexicon page: \url{https://cocoon.huma-num.fr/exist/crdo/lexvo.xql?code=uve}
\item A. Djoupa, doctoral dissertation on the syntax and semantics of Fagauvea (INALCO, 2013), cited on the same page.
\item Wikipedia overview: \url{https://en.wikipedia.org/wiki/West_Uvean_language}
\item Grambank entry: \url{https://grambank.clld.org/languages/west2516}
\end{itemize}

\section{Religious and Community Resources}
None found.

\section{Digital and Audio Resources}
\begin{itemize}
\item Omniglot page on the writing system: \url{https://www.omniglot.com/writing/westuvean.htm}
\item Endangered Languages Project entry: \url{https://endangeredlanguages.com/lang/11042}
\end{itemize}

\section{Curriculum Development Context}
Fagauvea has a dictionary and a doctoral grammar, but no teaching materials turned up. It is not known to be taught in schools. Any curriculum could draw on the Hollyman dictionary and the Djoupa grammar, once they are confirmed to be accessible, and should be developed with community input.
\end{document}